%% SYNTHETIC TEST UNIT generated by writing/tools/make_outline_book.py from FINAL_OUTLINE_v2 -- not book text
\chapter{矩问题与解析延拓引理}\label{app:appK}
本单元为测试用的合成文本（由 FINAL\_OUTLINE\_v2 自动生成），陈述内容从略。

\section{第 1 节（测试）}\label{sec:appK-hausdorff}

\begin{lemma}[Hausdorff 唯一性]\label{lem:appK-hausdorff}
（测试用合成陈述：内容从略。）
% uses: none
\end{lemma}

\begin{theorem}[偶数矩决定函数]\label{thm:appK-even-moments}
（测试用合成陈述：内容从略。）
\uses{lem:appK-hausdorff}
\end{theorem}

\begin{corollary}[PDF 的唯一性]\label{cor:appK-pdf}
（测试用合成陈述：内容从略。）
\uses{thm:appK-even-moments}
\end{corollary}

\section{第 2 节（测试）}\label{sec:appK-continuation}

\begin{proposition}[领头阶核由偶数矩唯一确定]\label{prop:appK-kernels}
（测试用合成陈述：内容从略。）
\uses{thm:appK-even-moments}
\end{proposition}

\begin{theorem}[Carlson 定理]\label{thm:appK-carlson}
（测试用合成陈述：内容从略。）
% uses: none
\end{theorem}

\begin{proposition}[多项式有界条件下延拓的唯一性]\label{prop:appK-continuation}
（测试用合成陈述：内容从略。）
% uses: none
\end{proposition}

\section{第 3 节（测试）}\label{sec:appK-entire}

\begin{theorem}[Müntz–Szász 定理]\label{thm:appK-muntz}
（测试用合成陈述：内容从略。）
% uses: none
\end{theorem}

\begin{lemma}[受控增长的整函数]\label{lem:appK-entire}
（测试用合成陈述：内容从略。）
% uses: none
\end{lemma}

\begin{lemma}[V. A. Markov 系数不等式]\label{lem:appK-markov}
（测试用合成陈述：内容从略。）
% uses: none
\end{lemma}

\section{第 4 节（测试）}\label{sec:appK-counterexamples}

\begin{remark}[反例与稳定性]\label{rem:appK-counterexamples}
（测试用合成陈述：内容从略。）
% uses: none
\end{remark}
